\section{Capítulo~\ref{sec:gaba}. \nt{GLUT} y \nt{GABA} juntos}

Las proposiciones lógicas y dinámicas del capítulo se deducen en él mismo a partir del análisis lineal y de la ley de Ohm; los experimentos muestran que los circuitos en los que se apoyan (el veto con retardo, la inhibición directa tras la excitación, la estabilización por inhibición, el ritmo del lazo \nt{GLUT}--\nt{GABA}) funcionan realmente en el cerebro.

{\footnotesize
\setlength{\tabcolsep}{3pt}\renewcommand{\arraystretch}{1.15}
\begin{longtable}{>{\raggedright}p{0.5cm}>{\raggedright}p{4.0cm}>{\raggedright}p{5.6cm}>{\raggedright}p{2.3cm}>{\raggedright\arraybackslash}p{1.7cm}}
\toprule
N.º & Afirmación & Experimento y qué se midió & Fuente & Estado \\
\midrule\endfirsthead
\toprule
N.º & Afirmación & Experimento y qué se midió & Fuente & Estado \\
\midrule\endhead
1 & Las neuronas \nt{GABA} son entre una quinta y una cuarta parte de las de la corteza & Inmunotinción de GABA en $10$ áreas de la corteza de mono: cerca del $25\%$ de las neuronas en $8$ áreas, menos del $20\%$ en la corteza visual primaria. Revisión de la diversidad de las neuronas inhibidoras de la corteza & \cite{Hendry1987,Markram2004} & medido \\
2 & Lo que hace la inhibición depende en gran medida del lugar de llegada: dendrita, cuerpo, lugar de nacimiento de la espiga, terminal de otro cable & Revisión: las clases de neuronas \nt{GABA} de la corteza se distinguen por su diana en el destinatario. Registro simultáneo del cuerpo y de una dendrita de una neurona piramidal: la inhibición directa es mucho más fuerte en el cuerpo que en las dendritas. La inhibición sobre la terminal de otro cable en la médula espinal reduce la liberación de neurotransmisor de una sola línea de entrada (revisión de mediciones) & \cite{Markram2004,Pouille2001,Rudomin1999} & medido \\
3 & El cambio de signo a través de un intermediario cuesta un retardo, cerca de un milisegundo; la inhibición que llega tras la excitación estrecha la ventana de coincidencia & Neuronas piramidales de hipocampo de rata en rebanadas: la inhibición a través de un intermediario llega tras la excitación directa con un retardo corto, y la ventana en la que dos entradas excitadoras se suman hasta el umbral es menor de $2\ms$. En las dendritas, donde esta inhibición es más débil, la ventana es más ancha & \cite{Pouille2001} & medido \\
4 & El veto $a\wedge\bar b$~\eqref{eq:veto} con OR y una constante uno es funcionalmente completo (proposición~\ref{prop:gabacomplete}) & Demostración en el capítulo. Resultado clásico: una red de elementos de umbral con entradas inhibidoras realiza cualquier expresión lógica. Proposición sobre el modelo; no hay experimento & \cite{McCullochPitts1943} & teoría \\
5 & El veto con retardo da un detector de dirección del movimiento con velocidad óptima $v^*=\Delta x/d$~\eqref{eq:vstar} & Retina de conejo: la selectividad a la dirección la crea una inhibición que, en el movimiento en la dirección “nula”, apaga la excitación. El registro por separado de las entradas excitadora e inhibidora de las células selectivas mostró que las células amacrinas estrelladas (\nt{GABA}) las inhiben de forma asimétrica, solo con las prolongaciones dirigidas hacia el lado “nulo”. La fórmula de $v^*$ es una deducción del capítulo & \cite{Barlow1965,Fried2002} & medido \\
6 & La inhibición de derivación divide el voltaje, no le resta~\eqref{eq:shunt} (fig.~\ref{fig:shunt}) & Ley de Ohm para un divisor de tres conductancias con el potencial de equilibrio del cloro cerca del reposo; para una neurona sin espigas & deducción del capítulo & teoría \\
7 & Sobre la frecuencia de espigas, la derivación actúa restando, y la división surge de la inhibición junto con un ruido de fondo equilibrado & Modelo: en una neurona que se dispara, la corriente por la derivación es casi constante y el efecto sobre la frecuencia es de resta. Experimento: en neuronas piramidales de corteza somatosensorial de rata (rebanadas) se inyectaron mediante pinza dinámica conductancias de fondo excitadoras e inhibidoras; el aumento simultáneo y equilibrado de ambas divide la pendiente de la respuesta sin un desplazamiento apreciable de la frecuencia. Revisión: la división por la actividad total se encuentra en muchas regiones del cerebro & \cite{HoltKoch1997,Chance2002,Carandini2012} & medido \\
8 & Con $\beta>1$, la inhibición mutua elige un solo ganador (proposición~\ref{prop:wta}, \eqref{eq:wta}) & Los valores propios $-(1\pm\beta)/\tau$ son una deducción del capítulo. Las frecuencias $0.73$ y $0.53$ con $\beta=0.5$ son el punto fijo $r_{1,2}=(I_{1,2}-\beta I_{2,1})/(1-\beta^2)$; no hay script en \texttt{models/}. El capítulo no cita un experimento directo & deducción del capítulo & teoría \\
9 & Reinicio de la memoria con una señal a través de \nt{GABA}; dos grupos con inhibición mutua forman un cerrojo & Se sigue de la fig.~\ref{fig:bistable} y de la proposición~\ref{prop:wta}. No hay experimento & deducción del capítulo & teoría \\
10 & El equilibrio del lazo \nt{GLUT}--\nt{GABA} es estable si la inhibición es suficientemente fuerte y suficientemente rápida (proposición~\ref{prop:ei}) & Modelo de dos poblaciones~\eqref{eq:ei}; las condiciones (i) y (ii) son el determinante y la traza de su matriz, deducidas en el capítulo & \cite{WilsonCowan1972} & teoría \\
11 & Con $w_{EE}=1.5$ y $w_{EI}w_{IE}=2$, la inhibición rápida ($\tau_I=2\ms$) mantiene $E^*=0.67h$, y la lenta ($30\ms$) desencadena oscilaciones (fig.~\ref{fig:ei}) & Solución numérica de~\eqref{eq:ei}, script \texttt{figures/make\_ei.py} & cálculo & modelo \\
12 & La corteza funciona en régimen de estabilización por inhibición; su firma: si se empuja a las neuronas \nt{GABA}, su frecuencia baja~\eqref{eq:isn} & La teoría predijo la respuesta paradójica. Experimento: registros intracelulares en la corteza visual primaria del gato; al suprimir la respuesta con un estímulo circundante, la conductancia inhibidora no crece, sino que cae junto con la excitadora. La excitación optogenética de neuronas PV del ratón en las cortezas visual, somatosensorial y motora reduce la frecuencia de las neuronas inhibidoras, también sin estímulo y con anestesia ligera. El coeficiente $-1/3$ con los números de la fig.~\ref{fig:ei} es una deducción del capítulo & \cite{Tsodyks1997isn,Ozeki2009,Sanzeni2020} & medido \\
13 & El lazo “\nt{GLUT} excita a \nt{GABA}, \nt{GABA} inhibe a \nt{GLUT}” genera un ritmo rápido con período de $12$--$50\ms$ & La estimulación optogenética de neuronas \nt{GABA} rápidas en la corteza somatosensorial del ratón a frecuencias de $8$--$200$\,Hz refuerza selectivamente el ritmo gamma ($20$--$80$\,Hz, período de $12$--$50\ms$); la estimulación de neuronas \nt{GLUT} refuerza solo ritmos más lentos & \cite{Cardin2009,Buzsaki2012} & medido \\
\bottomrule
\end{longtable}}
